% Verbatim copy of /Users/alingalatan/workspaces/grothe/handoff/manuscript/legacy/ch25-parts/40-regional-and-cumulant-caps.tex (file dated 2026-10-25). Only this line was added.
\subsection{From a joint support budget to a complete regional row cap}
\label{sec:canonical-response-cap}
\paragraph{Provenance (IDE note, 21 October).} Lower-bound GPT regional-cap package (\code{conjK-original-response-regional-cap-20261021.zip}, archived in
\pathref{gpt-reviews/regional-response-cap-20261021/}). IDE replays \code{python -I verify\_all.py} and \code{--fine} (pinned \code{python-flint 0.9.0}, unrelated directory):
ALL CLAIMS PASS (3.7\,s and 7.0\,s): 60 Young/optimal-$\eta$ audits, 16 simultaneous completions, the regional threshold, the larger quadratic region and the nonemptiness
witness. IDE re-derivation \st{P} of \eqref{eq:crc-anchor}: with $f=h+k$, $g=h-k$, $\mathcal N=t^2-\|Pk\|^2-\varepsilon\|Qh\|^2+p\|F\|^2$; then
$-\varepsilon\|Qh\|^2\le2\operatorname{Re}\langle B_0,h\rangle+2p\beta^2$; the support test with $L\pm q$ gives $2\ell t+2\operatorname{Re}\langle B_0,h\rangle+2p\|q\|^2\le\E|L+q|+\E|L-q|$
(using $\langle q,k\rangle=p\|q\|^2$); and $\|F\|^2-\|q\|^2=\|F-q\|^2$. So every visible sector enters only through $p\|F-q\|^2$, which the joint budget pays. The
cap inherits the remote theorem $\Lambda(q)<C'$ (Lean moment constants). It is a regional original-row cap: no stationarity, buffer, or raw-K statement is involved.


\paragraph{Status and scope.}
[P] denotes an analytic proof, [C] a directed or exact computation,
[F] a floating experiment, and [?] an outstanding obligation.
The new conclusion is a \emph{full regional bound on the original norm
row}, including all visible sectors and arbitrary orthogonal remote
components. It uses the previously certified remote cap with its stated
moment input. It proves neither raw Conjecture K nor a new global endpoint.
The region is nonempty: explicit unimodular responses are certified below.
No stationarity, independent residual buffer, or independence among
remote fields is assumed.

Fix $\varepsilon=1891/2500$, $p=1/\varepsilon$ and
\[
 h=(f_++f_-)/2,\quad k=(f_+-f_-)/2,\quad |f_\pm|\le1,
 \quad Ph=tZ,\quad F=\varepsilon Qk.
\]
Here $P,Q$ project onto first and $(2,1)$ chaos. All inner products use
$\langle u,v\rangle=\mathbb E[\bar u v]$; optimization uses their real parts.
Write $q=\mathbb E[F\mid X]$ for the remote cubic and use the decomposition
and $r_\perp,r_\parallel,E_{\rm ev}$ from the preceding joint-budget section.
Thus
\[
 \|F-q\|^2=\|r_\perp\|^2+\|r_\parallel\|^2+E_{\rm ev}.
\]
The exact original objective is
\begin{equation}\label{eq:crc-original}
 \mathcal N=\operatorname{Re}\langle f_+,(P-\varepsilon Q)f_-\rangle
 =t^2+p\|F\|^2-\|Pk\|^2-\varepsilon\|Qh\|^2.
\end{equation}
Cross terms cancel in the real part. This identity explains why an
arbitrary-cubic raw-K bound is stronger than the application needs.

\subsubsection{Canonical response dual: the auxiliary profile cancels [P]}
For $u,d\in\operatorname{ran}P$, $v,C\in\operatorname{ran}Q$, put
\[
 M(u,v,C,d)=\mathbb E|u+v+C+d|+\mathbb E|u+v-C-d|
                 +\|u\|^2+p\|C\|^2.
\]
Every actual response satisfies
\begin{equation}\label{eq:crc-canonical}
 \boxed{\mathcal N\le t^2+p\|F\|^2+M(u,v,C,d)
       +2p\operatorname{Re}\langle v,F\rangle
       +2\operatorname{Re}\langle d,tZ\rangle.}
\end{equation}
Indeed, majorize the two negative squares in \eqref{eq:crc-original} by
$\|u\|^2-2\operatorname{Re}\langle u,k\rangle$ and
$p\|C\|^2-2\operatorname{Re}\langle C,h\rangle$, and add the zero terms
$2\operatorname{Re}\langle v,pF-Qk\rangle$ and
$2\operatorname{Re}\langle d,tZ-Ph\rangle$.
Pointwise maximization over the two disks gives the two moduli.

For a fixed trial, $M$ is independent of the target $(t,F)$. Hence its
certified upper enclosure gives one quadratic cap reusable across rows.
In the previous four-block dual, with $L=\ell Z+B$, the substitutions
\[
 v=v_{\rm old}-F,\quad C=C_{\rm old}-B,
 \quad d=d_{\rm old}-\ell Z,
\]
followed by addition of the old row constant
$t^2-2\ell t+p\|B\|^2$, give \eqref{eq:crc-canonical} exactly.
This is a reparameterization of the retained-response dual, not a gain
obtained merely by changing a profile.

For completeness, in fixed finite dimension define the extended convex
function
\[
 J(t,F)=\inf\{\|Pk\|^2+\varepsilon\|Qh\|^2:
             |h\pm k|\le1,\ Ph=tZ,\ Qk=pF\}.
\]
It is closed: the disk-response set is weak-* compact and its finite
moment image is compact. Convexity follows from convexity of the squares.
The closed-value/Fenchel argument of the retained-penalty duality section,
applied to both equality constraints and both squares, shows that the
infimum of the right side of \eqref{eq:crc-canonical} equals
$t^2+p\|F\|^2-J(t,F)$. It is $-\infty$ at infeasible moments and need not
be attained on the boundary. Purification, if required, preserves the
finite $P,Q$ moments exactly. Only the upper-bound direction is needed here.

\subsubsection{A quadratic representation of the full parallel envelope [P]}
Use the previously proved notation
\[
 R=\sqrt\pi/2,\quad L_0=R\sqrt{3/8},\quad
 r(v)=2v\rho'(|v|^2),\quad
 D(v)=\mathbb E\rho(|v|^2)-t-p\operatorname{Re}\langle v,q\rangle,
\]
\[
 a(v)=-\frac R2P_X\{vG(|v|^2)\},\qquad
 G(u)=\frac2\pi\int_0^1e^{-us}\sqrt{s/(1-s)}\,ds.
\]
Here $v$ is any fixed remote cubic test. The previous parallel theorem is
$\Phi(M,B)+C_0\le D(v)$, where
\[
 M=\|\mathbb E[k\mid X]-r(v)\|,\quad
 B=\|p r_\parallel-\sqrt{3/2}\,a(v)\|,\quad
 C_0=p^2(\|r_\perp\|^2+E_{\rm ev})/4.
\]
It has the following exact convex-quadratic representation:
\begin{equation}\label{eq:crc-eta}
 \boxed{\Phi(M,B)=\sup_{0<\eta\le1}
 \left\{\frac{M^2}{2R+L_0^2(\eta^{-1}-1)}
                    +\frac{1-\eta}{4}B^2\right\}.}
\end{equation}
The endpoint $\eta=0$ means the limit $B^2/4$.
To prove it, use
\[
 (B-L_0s)_+^2=
 \sup_{0<\eta\le1}\{(1-\eta)B^2-(\eta^{-1}-1)L_0^2s^2\}
\]
in the defining supremum
$\Phi=\sup_{s\ge0}\{Ms-Rs^2/2+(B-L_0s)_+^2/4\}$.
The two suprema commute; completing the $s$ square proves the identity.
Thus the full nonquadratic envelope is exactly an intersection of
quadratic constraints. Keeping only finitely many $\eta$ weakens it safely.

For a finite remote cubic subspace $\mathcal B$, Bessel gives
\[
 M^2\ge \|P_Xk-P_Xr(v)\|^2
        +\|\Pi_{\mathcal B}(pq-r(v))\|^2.
\]
The exact inequalities $\sqrt\pi<1773/1000$ and
$L_0^2=3\pi/32<295/1000$, with $\eta=1/2$, yield
\begin{equation}\label{eq:crc-rational-cut}
\boxed{\begin{aligned}
 &\theta\left(\|P_Xk-P_Xr(v)\|^2+
                \|\Pi_{\mathcal B}(pq-r(v))\|^2\right)\\
 &\quad+\frac18\|p r_\parallel-\sqrt{3/2}\,a(v)\|^2
       +\frac{p^2}{4}(\|r_\perp\|^2+E_{\rm ev})\le D(v),
 \qquad \theta=\frac{250}{517}.
\end{aligned}}
\end{equation}
For a fixed test all centers are constants, so this is convex quadratic
in the actual response moments. It is not imposed on arbitrary raw-K cubics.

\subsubsection{Remote anchoring and joint payment of every visible sector [P]}
Put $\ell=1057/1250$, $\beta=179/1250$,
$B_0=\beta Z(|Z|^2-2)$, and $L=\ell Z+B_0$.
Taking $u=0$, $v=-q$, $C=-B_0$, $d=-\ell Z$ in
\eqref{eq:crc-canonical}, but leaving the negative $\|Pk\|^2$ unrelaxed,
gives the useful anchor inequality
\begin{equation}\label{eq:crc-anchor}
 \boxed{\mathcal N\le t^2-2\ell t+2p\beta^2+\Lambda(q)
                  +p\|F-q\|^2-\|Pk\|^2.}
\end{equation}
This can also be proved directly from
$2\ell t+2\operatorname{Re}\langle B_0,Qh\rangle+2p\|q\|^2
\le\mathbb E|L+q|+\mathbb E|L-q|$ and completion of the $Qh$ square.
It retains arbitrary interactions among every visible sector.

Assume the inherited remote theorem $\Lambda(q)<C'$ with
$C'=2747927/2000000$. Its hypotheses are the project's radial moment
bounds, including $C_6=4045313/100000$; this package does not reprove them
or regenerate that remote certificate. For any $\mu>8/p$, let
\[
 m^2=\|\Pi_{\mathcal B}(pq-r(v))\|^2,\quad
 a_1=P_Xr(v),\quad a_2=\sqrt{3/2}\,a(v),\quad
 E_0=\|r_\perp\|^2+E_{\rm ev}.
\]
Combining \eqref{eq:crc-rational-cut} and \eqref{eq:crc-anchor} proves
\begin{equation}\label{eq:crc-transfer}
\boxed{\begin{aligned}
\mathcal N <{}& t^2-2\ell t+2p\beta^2+C'
 +\mu\{D(v)-\theta m^2\}
 -\frac{\mu\theta}{1+\mu\theta}\|a_1\|^2
 +\frac{\mu}{\mu p-8}\|a_2\|^2\\
 &-\left(\frac{\mu p^2}{4}-p\right)E_0 .
\end{aligned}}
\end{equation}
No maximizer assumption is used. In particular this is different from
the stationarity-only Ward estimate.

\paragraph{The simultaneous completion.}
Write $x=P_Xk$ and $y=p r_\parallel$. After multiplying the joint cut by
$\mu$, the two remaining suprema are
\[
 \sup_x\{-\|x\|^2-\mu\theta\|x-a_1\|^2\}
       =-\frac{\mu\theta}{1+\mu\theta}\|a_1\|^2,
\]
\[
 \sup_y\{\|y\|^2/p-\mu\|y-a_2\|^2/8\}
       =\frac{\mu}{\mu p-8}\|a_2\|^2.
\]
The other sector coefficient is $p-\mu p^2/4<0$.
These are joint completions of one valid constraint; separate optimized
penalties have not been added. The result is already a cap on the full
original objective and is not an additional deduction from an optimized
four-block response dual.

\subsubsection{A nonempty, dimension-free regional row certificate [P,C]}
Choose any nine orthonormal remote complex Gaussian directions
$X,Y_1,\ldots,Y_4,W_1,\ldots,W_4$ and set
\[
 H=X(|X|^2-2)/\sqrt2,\qquad
 J=\frac X2\sum_{j=1}^4Y_j\bar W_j,\qquad
 V_0=H+\frac{113}{40}J.
\]
$H,J$ are orthonormal cubics. All extra remote directions and all
orthogonal cubic coefficients are unrestricted. The theorem is invariant
under changing this orthonormal frame; it therefore applies to the union
over all such frames, not only a single coordinate choice.

\paragraph{Regional theorem.}
If an actual response satisfies
\begin{equation}\label{eq:crc-region}
 t\ge\frac{1057}{1250},\qquad
 \operatorname{Re}\langle V_0,q\rangle\ge\frac{58969}{100000},
\end{equation}
then
\begin{equation}\label{eq:crc-cap}
 \boxed{\mathcal N<\frac{713260}{10^6}
       -5(\|r_\perp\|^2+E_{\rm ev})<\frac{500}{701}.}
\end{equation}
The last strict inequality holds also when the extra sector energy is zero.
This includes arbitrary shared $A/B$, even and diagonal sectors, at every
amplitude compatible with an actual response. No buffer is introduced.
The inherited remote moment hypothesis remains explicit.

\paragraph{A larger rational quadratic region.}
The same conclusion holds if the linear condition in
\eqref{eq:crc-region} is replaced by
\begin{equation}\label{eq:crc-quadratic-region}
 p s\operatorname{Re}\langle V_0,q\rangle
 +\theta\left(\left|p\langle H,q\rangle-\frac{8027557}{10^8}\right|^2
             +\left|p\langle J,q\rangle-\frac{24750622}{10^8}\right|^2\right)
 \ge p s\frac{58969}{100000},\qquad s=\frac{221}{2000}.
\end{equation}
This contains the entire linear region. The rounded cubic profile center
has norm error below $3\cdot10^{-9}$. Since $\|pq\|\le1$ for an actual
response and the rounded center has norm below $0.27$, replacing the
true squared projection defect by this rational one costs at most
$10^{-8}$. Multiplication by $\mu\theta$ adds less than
$7.02\cdot10^{-8}$ to the certified cap, still below $0.713260$.
This uses the nonpositive projection term in \eqref{eq:crc-transfer},
not an additional moment assumption.

\paragraph{Certified constants.}
Use $v=(221/2000)V_0$ and $\mu=29/2$ in
\eqref{eq:crc-transfer}, and safely drop the nonpositive $-\mu\theta m^2$.
Phase symmetry gives $a_1=\xi X$ and $a(v)=\alpha X$ with real scalars.
Fresh directed integration proves
\[
\begin{split}
 \mathbb E\rho(|v|^2)&=0.931750156440\ldots,\\
 \xi&=-0.005846273982\ldots,\qquad
 \alpha=0.006507965772\ldots,\\
 D(v)&\le 0.000004274632\ldots\quad\hbox{under \eqref{eq:crc-region} at }t=\ell.
\end{split}
\]
The interval calculation of the complete right side is
\[
 \mathcal N+5E_0 <0.713259308<0.713260,
\]
with the default integrator error below $2\cdot10^{-9}$ for that bound.
The exact coefficient $\mu p^2/4-p$ exceeds $5$.
For larger $t$, the right side changes by
$(t-\ell)^2-\mu(t-\ell)\le0$, because every actual response has $t\le1$.
The result therefore holds throughout $t\ge\ell$, not only at one node.

\paragraph{The region contains actual responses.}
Let $s_*=55263/500000$ and take
\[
 f_\pm=\operatorname{phase}(Z\pm s_*V_0).
\]
Then $Ph=t_*Z$, and for $F_*=\varepsilon Qk_*$,
$q_*=\mathbb E[F_*\mid X,Y,W]$, the fresh verifier proves
\[
 t_*=0.8456011221357\ldots>0.8456,\qquad
 \operatorname{Re}\langle V_0,q_*\rangle
     =0.5897115663962\ldots>0.58969.
\]
Thus \eqref{eq:crc-cap} is not merely an assertion about an empty slice.
The support identity gives these numbers exactly from
\[
 t_* =\mathbb E\rho(|s_*V_0|^2)
       -s_*\operatorname{Re}\langle V_0,r(s_*V_0)\rangle,
 \qquad
 \operatorname{Re}\langle V_0,q_*\rangle
       =\varepsilon\operatorname{Re}\langle V_0,r(s_*V_0)\rangle.
\]
Visible parity kills the other components of $Ph$.
This establishes membership; it does not claim that these responses
maximize the norm or that the new region contains the global bottleneck.

\subsubsection{Directed integral formulas and tails [C]}
Condition on $T=|X|^2$ and $Y=\sum_{j=1}^4|W_j|^2$.
They have independent densities $e^{-T}$ and $Y^3e^{-Y}/6$.
For a test $v=cH+dJ$ with real coefficients used here, put
\[
 Z_T=T(T-2)^2/2,\quad S_T=1+d^2TY/4,\quad
 u=c^2Z_T/S_T,\quad B_j=e^{-u/2}I_j(u/2).
\]
The needed integrands, before multiplying by the two densities, are
\[
\begin{array}{ll}
 \mathbb E\rho(|v|^2):&R\sqrt{S_T}\{(1+u)B_0+uB_1\},\\
 \langle X,r(v)\rangle:&Rc\,T(T-2)(B_0+B_1)/\sqrt{2S_T},\\
 \langle H,r(v)\rangle:&Rc\,Z_T(B_0+B_1)/\sqrt{S_T},\\
 \langle J,r(v)\rangle:&Rd\,TYB_0/(4\sqrt{S_T}),\\
 \langle X,a(v)\rangle:&-Rc\,T(T-2)(B_0-B_1)/(2\sqrt2 S_T^{3/2}).
\end{array}
\]
The last formula uses the Gaussian convolution identity
$\mathbb E[(w+\sqrt h U)G(|w+\sqrt h U|^2)]
=w(1+h)^{-3/2}G(|w|^2/(1+h))$.
It follows by writing $G$ as its positive Laplace integral and making
$s\mapsto s/(1+hs)$, or by Gaussian differentiation.

The integration rectangle is $[0,36]\times[0,48]$.
For the first four integrands, the nonnegative polynomial tail envelopes
are respectively
\[
 3+4T+T^2+Y/4,\quad T^2+2T,\quad
 T^3/2+2T^2+2T,\quad TY/4.
\]
The new profile envelope is $T^2/2+T$.
These apply because $cR,dR,c<1$ in both evaluations, $S_T\ge1$,
and the positive Gaussian kernels are bounded by one where used.
Integrating the polynomials over the union $T>36$ or $Y>48$ gives explicit
incomplete-gamma tails. Complex-domain bounds use $\operatorname{Re}S_T>0$
and the positive Laplace kernels, so Arb verifies holomorphy instead of
silently treating real quadrature nodes as a certificate.

\subsubsection{A finite quadratic certificate interface [P]}
For real moment coordinates $x$, every fixed canonical trial has form
$U_i(x)=x^\top Q_0x+2b_i^\top x+a_i$.
Necessary quadratic constraints have form
$g_j(x)=x^\top G_jx+2h_j^\top x+k_j\le0$.
If $\theta_i,\lambda_j\ge0$, $\sum_i\theta_i=1$, and
\[
 \begin{pmatrix}
 \sum_j\lambda_jG_j-Q_0&\sum_j\lambda_jh_j-\sum_i\theta_i b_i\\
 *&B-\sum_i\theta_i a_i+\sum_j\lambda_j k_j
 \end{pmatrix}\succeq0,
\]
then $\mathcal N\le\min_iU_i(x)\le B$ on that feasible set.
The matrix says $B-\sum_i\theta_iU_i+\sum_j\lambda_jg_j\ge0$;
feasibility and nonnegative multipliers prove the claim.
A rational LDL certificate suffices once all integral constants are safely
enclosed. For instance, an approximate center with error $\delta$ obeys
$\|z-r\|^2\ge(1-s)\|z-\hat r\|^2-(s^{-1}-1)\delta^2$ for $0<s<1$.
Use upper enclosures for the support integral and the cap constants.
No successful global numerical matrix certificate is claimed in this package.

\subsubsection{What the numerical implementation did and did not establish [F,?]}
The floating code optimizes all twelve complex visible cubic coefficients
on $(Z,X_1,X_2)$, a separate independent grade-three Gaussian limit
coefficient, and fixed remote tensor shapes, while keeping the actual
$t=0.8456$. Fixed dual trials are re-evaluated on independent grids.
Some candidate deductions are several $10^{-5}$, but near the Gaussian
benchmark they are below $10^{-6}$. Removing the joint support terms
changes the final local outer values by only $5\cdot10^{-10}$ to
$1.3\cdot10^{-6}$, below the unresolved optimization gaps in the larger
cases. It does not demonstrate a robust gain attributable to these cuts
on the hardest tensors. A whole-ray semidefinite relaxation was also too
loose and had unresolved numerical residuals. These runs are diagnostics,
not certificates or a row cover.

The certified progress is \eqref{eq:crc-cap}: an entire nonempty region of
the \emph{full original row} is now below the desired cap. The remaining
obligation is to bound its complement, or construct a library of certified
support-face regions sufficient for every potentially extremizing tensor,
and then cover all rows at the same $\varepsilon$. Dimension-free coverage
is not supplied by a finite coordinate search. The certified global
Grothendieck bounds therefore remain unchanged.

\emph{Update (dimension-free criterion):} \S\ref{sec:cumulant-response-cap} gives a dimension-free sufficient condition for the central row cap
($\kappa_v\le10^{-3}$, $D\le5\cdot10^{-6}$ for one remote test, e.g.\ $v=(p/t)q$). The obligation to cover the complement of all certified regions remains; the result
is regional, not global. The complementary region with actual remote variance at most $0.0148$ is covered throughout the central interval by \S\ref{sec:anisotropic-response-cap},
with all shared sectors and no moment hypothesis. \S\ref{sec:response-interior-operator} remains available on its own slightly wider $t$ interval. The intermediate-variance
and larger-cumulant/influence complement remains open.

\subsection{A cumulant--support criterion for the full original row}
\label{sec:cumulant-response-cap}
\paragraph{Provenance (IDE note, 21 October).} Lower-bound GPT cumulant-response package (\code{conjK-cumulant-response-cap-20261021.zip}, archived in
\pathref{gpt-reviews/cumulant-response-cap-20261021/}). IDE replay \code{python -I verify\_all.py} (unrelated directory): ALL CLAIMS PASS (39 exact checks). IDE check
\st{P} of the assembly: the anchor inequality is that of \S\ref{sec:canonical-response-cap} (IDE re-derived there); the support inequality \eqref{eq:cr-support-vector} places
the A/B/even coefficient vector within $2\sqrt D$ of a centre of norm $\sqrt{3/2}\|a(v)\|\le\chi$, so $p\|F-q\|\le2\sqrt D+\chi$ and $p\|F-q\|^2\le\varepsilon(2\sqrt D+\chi)^2$.
The analytic input \eqref{eq:cr-profile} uses the resolvent lemma, H1 and the influence estimate; the replay certifies only arithmetic. The remote cap $C'$ keeps its
manuscript moment hypotheses. This criterion concerns original responses, not raw K; the warning about steep $L$-profile derivatives still stands.


\paragraph{Scope and status.}
The results in this section are [P] analytic proofs and [C] exact rational
checks. They concern actual bounded responses in the original norm problem,
not arbitrary perturbations in raw Conjecture K. They hold in every finite
dimension, with all shared sectors present, all complex phases, and no
independent Gaussian buffer. The analytic inputs are the established H1
fluctuation inequality, the influence inequality, and the original-response
support constraint. Numerical row conclusions additionally use the existing
remote cap with its stated moment input. No global endpoint is claimed.

Put
\[
 \varepsilon=\frac{1891}{2500},\quad p=\frac1\varepsilon,
 \quad \ell=\frac{1057}{1250},\quad\beta=\frac{179}{1250},
 \quad L=\ell Z+\beta Z(|Z|^2-2).
\]
Let $P=P_{1,0}$ and $Q=P_{2,1}$. For disc-valued responses $f_+,f_-$,
write
\[
 h=(f_++f_-)/2,\quad k=(f_+-f_-)/2,\quad Ph=tZ,
 \quad F=\varepsilon Qk,\quad q=\mathbb E[F\mid X].
\]
Here $Z$ is the visible standard proper complex Gaussian, independent of
the remote Gaussian vector $X$. The original objective is
\[
 \mathcal N=\operatorname{Re}\langle f_+,(P-\varepsilon Q)f_-\rangle
 =t^2+p\|F\|_2^2-\|Pk\|_2^2-\varepsilon\|Qh\|_2^2.
\]
The usual visible decomposition is
\[
 F=q+(|Z|^2-1)r_A+Z^2r_B/\sqrt2+ZY+\bar ZW+\nu H(Z),
 \qquad H(Z)=Z(|Z|^2-2)/\sqrt2.
\]
No independence among the remote fields is imposed.

For a remote cubic test $v\in H_{2,1}(X)$, define
\begin{align*}
 \kappa_v&=\mathbb E|v|^4-2\|v\|_2^4,\\
 I_v&=\sup_{\|b\|=1}
   \left(\|D_{\bar b}v\|_2^2+\|\bar D_bv\|_2^2\right),\\
 D_t(v;q)&=\mathbb E\rho(|v|^2)-t-p\operatorname{Re}\langle v,q\rangle,
 \qquad \rho(u)=\mathbb E|Z+\sqrt u|,\\
 \chi(v)&=\min\left\{\kappa_v^{3/4},
                   \sqrt{\frac87\kappa_v I_v}\right\}.
\end{align*}
For actual responses, $D_t(v;q)\ge0$ for every such test.
The test is free: it need not be proportional to $q$, nor use the same
influential coordinates as the visible sectors.

\subsubsection{The resulting complete-row criterion [P]}

\paragraph{Theorem.}
For every remote cubic test $v$,
\begin{equation}\label{eq:cr-main}
 \boxed{\quad
 \mathcal N\le t^2-2\ell t+2p\beta^2+\Lambda(q)
       +\varepsilon\left(2\sqrt{D_t(v;q)}+\chi(v)\right)^2,
 \quad}
\end{equation}
where
\[
 \Lambda(q)=\mathbb E(|L+q|+|L-q|)-p\|q\|_2^2.
\]
In particular the same bound holds after taking the infimum of the
nonnegative charge over all remote cubic tests. This is a universal
response inequality, not a dimension-dependent search prescription.

\paragraph{Certified central region.}
Assume the inherited remote theorem
\[
 \Lambda(q)<C':=\frac{2747927}{2000000}.
\]
Then
\begin{equation}\label{eq:cr-region}
 \boxed{\quad
 \frac{21}{25}\le t\le\frac{851}{1000},\qquad
 \kappa_v\le\frac1{1000},\qquad
 D_t(v;q)\le\frac1{200000}
 \quad\Longrightarrow\quad
 \mathcal N<\frac{713254}{10^6}<\frac{500}{701}.
 \quad}
\end{equation}
The assumption is existence of one test satisfying the two inequalities.
All visible sectors and all other remote cubic coefficients are allowed.
The cap is for the whole central interval, not only its peak.

The exact rational replay also checks the alternative sufficient pairs
\[
 (\kappa_v,D_t)\le(1/500,1/2000000)
 \quad\hbox{or}\quad
 (\kappa_v,D_t)\le(1/432,0),
\]
each giving a strict cap below $500/701$ on the same interval. These
alternatives use the same analytic theorem; the first displayed region
is the main numerical claim.

\subsubsection{A smooth profile in place of the stiff content profile [P]}

Put $R=\sqrt\pi/2$ and
\[
 G(u)=\frac2\pi\int_0^1e^{-us}\sqrt{\frac{s}{1-s}}\,ds,
 \qquad a(v)=-\frac R2P_X\{vG(|v|^2)\},
\]
where $P_X$ is remote first-chaos projection. The existing visible
Gaussian score calculation gives this exact $A/B$ profile in the
original-response support constraint. The following estimate is new:
\begin{equation}\label{eq:cr-profile}
 \boxed{\quad
 \frac32\|a(v)\|_2^2
 \le\min\left\{\kappa_v^{3/2},\frac87\kappa_v I_v\right\}.
 \quad}
\end{equation}
It treats $G$ exactly as a positive Laplace mixture; it is not a small
amplitude expansion. Its derivatives are on the unit-Gaussian scale,
unlike the much steeper nonlinear $L$-profiles in the older content
inequality. Thus the failure of those older sup-norm estimates is not
being ignored or reversed.

Here is a proof, including the numerical constants and the Gaussian
facts it uses.

\paragraph{Second-chaos resolvent lemma.}
Let $v\in H_{2,1}$, $e=\mathbb E|v|^2>0$, and let $\Pi$ be either
$\Pi_{1,1}$ or $\Pi_{2,0}$. If $\psi:\mathbb C\to\mathbb C$ has
charge zero or two, respectively, with bounded second derivatives,
then
\begin{equation}\label{eq:cr-resolvent}
 \|\Pi\psi(v)\|_2^2\le
 \frac3{16}\left(4\|\psi_{u\bar u}\|_\infty^2+
  2(\|\psi_{uu}\|_\infty+\|\psi_{\bar u\bar u}\|_\infty)^2\right)
 \kappa_v.
\end{equation}
It suffices here to use smooth functions of polynomial growth; all
profiles below satisfy these conditions. The zero polynomial is handled
separately.

To prove the lemma, let $\mathcal L$ be the real Gaussian
Ornstein--Uhlenbeck generator, normalized so that $\mathcal Lv=-3v$,
and let $\Gamma$ be its carr\'e du champ. Set
\[
 w=\Gamma(v,\bar v)/3,\qquad c=\Gamma(v,v)/3.
\]
The established H1 input is
\begin{equation}\label{eq:cr-h1}
 \mathbb Ew=e,\qquad
 \|w-e\|_2^2+\tfrac12\|c\|_2^2\le\kappa_v/3.
\end{equation}
On the scalar complex Gaussian space of variance $e$, use the positive
number operator
\[
 N_e=u\partial_u+\bar u\partial_{\bar u}
                   -2e\partial_u\partial_{\bar u}.
\]
After removing the Gaussian mean of $\psi$, its Hermite expansion has
no degrees zero or one: charge zero removes degree one, and charge two
removes both. Let $T_s=e^{-sN_e}$ and set
\[
 H(u)=\int_0^\infty e^{2s}T_{3s}
       (\psi-\mathbb E\psi(\sqrt e Z))(u)\,ds.
\]
Then $(3N_e-2)H=\psi-\mathbb E\psi(\sqrt eZ)$.
The spectral gap just stated gives convergence. Alternatively the
Mehler formula, subtraction of the mean and linear terms, and Taylor's
formula give decay $O(e^{-6s}(1+|u|^2))$.
Commutation of each scalar derivative with the Mehler operator gives
\[
 \|H_{uu}\|_\infty\le\tfrac14\|\psi_{uu}\|_\infty,
 \quad\|H_{u\bar u}\|_\infty\le\tfrac14\|\psi_{u\bar u}\|_\infty,
 \quad\|H_{\bar u\bar u}\|_\infty
                   \le\tfrac14\|\psi_{\bar u\bar u}\|_\infty.
\]
Indeed the factor in the integral is $e^{2s}e^{-6s}$.

The Gaussian chain rule is
\[
 \mathcal L H(v)=-3(N_eH)(v)+3\{cH_{uu}(v)
                +2(w-e)H_{u\bar u}(v)+\bar cH_{\bar u\bar u}(v)\}.
\]
Since $\Pi\mathcal L=-2\Pi$, projection gives
\[
 \Pi\psi(v)=3\Pi\{cH_{uu}(v)+2(w-e)H_{u\bar u}(v)
                                  +\bar cH_{\bar u\bar u}(v)\}.
\]
Projection is a contraction. Cauchy--Schwarz with
\eqref{eq:cr-h1} and the preceding derivative bounds proves
\eqref{eq:cr-resolvent}. The finite-dimensional Gaussian integration
by parts and polynomial growth justify these operations; no limiting
dimension or assumption on independent components is used.

\paragraph{From second-chaos to first-chaos content.}
For a charge-one function $F$ and a unit first-chaos direction $b$,
Gaussian integration by parts gives
\[
 \mathbb E[\bar bF(v)]
 =\mathbb E[F_u(v)D_{\bar b}v]
       +\mathbb E[F_{\bar u}(v)\overline{\bar D_bv}].
\]
The derivatives of $v$ are in $H_{1,1}$ and $H_{2,0}$. Consequently
\begin{equation}\label{eq:cr-content}
 \|P_XF(v)\|_2^2\le I_v
     \left(\|\Pi_{1,1}F_u(v)\|_2^2+
           \|\Pi_{2,0}F_{\bar u}(v)\|_2^2\right).
\end{equation}
The established influence contraction inequality supplies
\begin{equation}\label{eq:cr-influence}
 I_v\le\sqrt{3\kappa_v}/2.
\end{equation}
Keeping the actual $I_v$ can be considerably stronger than substituting
\eqref{eq:cr-influence}.

\paragraph{Global derivative bounds for $F(u)=uG(|u|^2)$.}
The probability measure defining $G$ has mean $3/4$. With $x\ge0$,
\begin{equation}\label{eq:cr-exp-bounds}
 |(-2+4x-x^2)e^{-x}|\le2,\qquad
 |x(3-x)e^{-x}|\le\tfrac45,\qquad
 x^2e^{-x}\le\tfrac{11}{20}.
\end{equation}
The first follows from $e^x\ge1+x+x^2/2$: the two residual
polynomials are $2(x-1/2)^2+7/2$ and $6x$.
For the second, the positive branch has its maximum at
$x_*=(5-\sqrt{13})/2\in[6972/10000,6973/10000]$.
The rational Taylor lower bound for $e^{6972/10000}$ through degree
twelve proves $x_*(3-x_*)e^{-x_*}<4/5$.
On $x\ge3$ use $x(x-3)e^{-x}\le x^2e^{-x}$.
The latter is at most $4e^{-2}<11/20$, checked with the degree-six
exponential polynomial. These are global bounds, including the tails.

For $u\in\mathbb C$ put $s=|u|^2$. Direct differentiation yields
\begin{align*}
 F_{uu\bar u}&=2G'(s)+4sG''(s)+s^2G'''(s),\\
 F_{uuu}&=\bar u^2\{3G''(s)+sG'''(s)\},\\
 F_{u\bar u\bar u}&=u^2\{3G''(s)+sG'''(s)\},\\
 F_{\bar u\bar u\bar u}&=u^4G'''(s).
\end{align*}
Integrating \eqref{eq:cr-exp-bounds} against the positive Laplace
measure gives
\[
 \|F_{uu\bar u}\|_\infty\le\frac32,
 \quad\|F_{uuu}\|_\infty,\ \|F_{u\bar u\bar u}\|_\infty\le\frac35,
 \quad\|F_{\bar u\bar u\bar u}\|_\infty\le\frac{33}{80}.
\]
Apply \eqref{eq:cr-resolvent} separately to $\psi=F_u$ and
$\psi=F_{\bar u}$. Adding the two squared estimates gives
\begin{align*}
 \|\Pi_{1,1}F_u(v)\|_2^2+\|\Pi_{2,0}F_{\bar u}(v)\|_2^2
 &\le C_2\kappa_v,\\
 C_2&=\frac3{16}\left\{4(3/2)^2+2(6/5)^2
                 +4(3/5)^2+2(3/2+33/80)^2\right\}
       =\frac{198099}{51200}.
\end{align*}
This sum explicitly includes both second-chaos terms.
Equations \eqref{eq:cr-content} and $a=-(\sqrt\pi/4)P_XF(v)$ imply
\[
 \|a(v)\|_2^2\le\frac\pi{16}C_2\kappa_v I_v.
\]
Using $\pi<22/7$ and $\sqrt3<7/4$ gives the exact comparisons
\[
 \frac{3\pi C_2}{32}<\frac{6537267}{5734400}<\frac87,
 \qquad
 \frac{\pi\sqrt3 C_2}{32}
       <\frac{2179089}{3276800}<\frac23.
\]
This proves \eqref{eq:cr-profile} in both forms. The proof does not
estimate the high derivatives of the original nonlinear $L$-profile.

\subsubsection{Paying for every visible sector [P]}

The established original-response support constraint is
\begin{equation}\label{eq:cr-support-vector}
 \|pr_A-a(v)\|_2^2+
 \|pr_B-\overline{a(v)}/\sqrt2\|_2^2
       +p^2(\|Y\|_2^2+\|W\|_2^2+|\nu|^2)
 \le4D_t(v;q).
\end{equation}
For completeness, it follows by testing boundedness with $v+U$, where
$U$ is an arbitrary visible-sector cubic, using the norm majorant,
averaging $Z$ first, and completing squares. The $A/B$ and even blocks
have uniform inverse-distance constant two. The even linear term and
the cross term vanish by $Z\mapsto-Z$. The $A/B$ linear term is
$(a(v),\overline{a(v)}/\sqrt2)$. Each assertion is conditional on the
entire remote vector, so shared dependence is preserved.

View all coefficients in \eqref{eq:cr-support-vector} as one Hilbert
vector. Its center has norm $\sqrt{3/2}\,\|a(v)\|_2\le\chi(v)$.
The triangle inequality gives
\begin{equation}\label{eq:cr-visible-energy}
 p\|F-q\|_2\le2\sqrt{D_t(v;q)}+\chi(v).
\end{equation}
This is one simultaneous budget for all visible sectors. It does not
add separately optimized costs.

The response dual anchored at the actual remote projection gives
\[
 \mathcal N\le t^2-2\ell t+2p\beta^2+\Lambda(q)
                  +p\|F-q\|_2^2-\|Pk\|_2^2.
\]
Indeed boundedness implies
\[
 2\ell t+2\operatorname{Re}\langle\beta Z(|Z|^2-2),Qh\rangle
                   +2p\|q\|_2^2
 \le\mathbb E(|L+q|+|L-q|),
\]
and completing the $Qh$ square proves the assertion.
Drop the nonpositive first-chaos square and apply
\eqref{eq:cr-visible-energy}. Since $1/p=\varepsilon$, this is
\eqref{eq:cr-main}.

\subsubsection{Exact numerical comparisons [C]}

On $t\in[21/25,851/1000]$,
\[
 t^2-2\ell t+2p\beta^2+C'
 \le\frac{2697232101}{3782000000}.
\]
At the hypotheses of \eqref{eq:cr-region}, exact root enclosures are
\[
 2\sqrt{D_t}\le\frac{447214}{10^8},\qquad
 \kappa_v^{3/4}\le\frac{562342}{10^8}.
\]
The checker verifies the appropriate second and fourth powers rather
than using floating square roots. Therefore the row cap is strictly below
\[
 \frac{2107440363267305224801}{2954687500000000000000}
 =0.713253216\ldots<0.713254<500/701.
\]
The gap between $0.713254$ and $500/701$ exceeds $1.27\cdot10^{-5}$.
This is a certified \emph{regional} row, not a new global lower bound.

\subsubsection{A canonical scalar test and a nonempty diffuse family [P,C]}

\paragraph{Collinear choice, without a tensor-frame search.}
If $t>0$, the choice $v=(p/t)q$ yields
\[
 D_t=\mathbb E\rho(p^2|q|^2/t^2)-t-\frac{p^2}{t}\|q\|_2^2,
 \quad\kappa_v=(p/t)^4\kappa_q,\quad I_v=(p/t)^2I_q.
\]
Thus one may use
\[
 \chi(v)=(p/t)^3\min\{\kappa_q^{3/4},\sqrt{8\kappa_qI_q/7}\}.
\]
This choice is the support optimizer for a proper Gaussian remote law
on its feasible boundary, but the displayed formulas and theorem are
valid for every remote cubic. A one-dimensional scale search
$v=\lambda q$ can improve it; noncollinear tests remain allowed.

\paragraph{An integration-free sufficient condition.}
The support gap can itself be bounded with three radial moments. The
unit Rice mean has the global cubic majorant
\begin{equation}\label{eq:cr-rice-cubic}
 \rho(u)\le R\left(1+\frac u2-\frac{u^2}{16}+\frac{u^3}{96}\right),
 \qquad u\ge0.
\end{equation}
Indeed $2\rho'(u)/R$ is the Laplace transform of the beta probability
law with parameters $(1/2,3/2)$. Its first two moments are $1/4$ and
$1/8$, giving the printed coefficients, and $\rho^{(4)}\le0$ gives
the sign of the Taylor remainder on the whole half-line.
If $R_+\ge R$ is a rational upper bound, set
\[
 \widehat D_t(v;q)=R_+\left(1+\frac{\mathbb E|v|^2}{2}
      -\frac{\mathbb E|v|^4}{16}+\frac{\mathbb E|v|^6}{96}\right)
       -t-p\operatorname{Re}\langle v,q\rangle.
\]
For actual responses $0\le D_t\le\widehat D_t$. The bracket is
positive because its integrand majorizes the positive Rice mean.
Thus $\widehat D_t<0$ proves infeasibility; otherwise it may replace
$D_t$ in every sufficient condition above. At $v=\lambda q$, this
requires only the second, fourth and sixth moments of $q$ and its
influence bound. The majorant is a sufficient test and can lose
accuracy; a sharper certified Rice integral remains permissible.

For $t\in[1057/1250,851/1000]$, the established response energy bound
$\|q\|^2<0.040726$ implies that normalized fourth moment at most $2.1$
gives $(p/t)^4\kappa_q<1/1000$. In that range only the displayed support
gap condition remains to apply \eqref{eq:cr-region}. This does not say
that every such law has small support gap.

\paragraph{Explicit responses showing nonemptiness.}
For independent standard proper complex $X_1,\ldots,X_N$, put
\[
 V_N=\frac3{10\sqrt N}\sum_{j=1}^N H(X_j),
 \qquad f_\pm=\operatorname{phase}(Z\pm V_N).
\]
Construct $h,k,F,q$ from these responses as above. Then
\begin{equation}\label{eq:cr-wick-moments}
 e_V=\frac9{100},\qquad
 \kappa_{V_N}=\frac{42e_V^2}{N},\qquad I_{V_N}=\frac{3e_V}{N}.
\end{equation}
Hence the influence form of the new profile bound gives
\[
 \chi(V_N)\le\sqrt{\frac87\kappa_{V_N}I_{V_N}}
              =\frac{81}{250N}.
\]
By equality in the defining support inequality,
\[
 D_t(V_N;q)=0.
\]
This is exact: $p\operatorname{Re}\langle V_N,q\rangle
=\mathbb E[2|V_N|^2\rho'(|V_N|^2)]$ and
$t=\mathbb E[\rho(|V_N|^2)-2|V_N|^2\rho'(|V_N|^2)]$.
The remote projection $q$ is the projection generated by these
responses; it is not being prescribed to equal a Wick sum.

The moment identities used to locate the row are
\[
 \mathbb E|V_N|^4=e_V^2(2+42/N),\qquad
 \mathbb E|V_N|^6=e_V^3(6+378/N+10968/N^2).
\]
They follow from $\mathbb E|H|^2=1$, $\mathbb E|H|^4=44$,
$\mathbb E|H|^6=11352$ and independence/circularity. In the fourth
moment the partition types are $2$ and $1+1$; in the sixth they are
$3$, $2+1$, and $1+1+1$, with multiplicities
$N$, $9N(N-1)$, and $36\binom N3$ before normalization.

To bound $t$ without quadrature, use
\[
 \rho(u)-2u\rho'(u)=R e^{-u/2}I_0(u/2)
 =\frac R\pi\int_0^\pi e^{-u(1-\cos\theta)/2}\,d\theta.
\]
The second and third exponential Taylor polynomials give
\[
 R\left(1-\frac{e_V}2+\frac3{16}\mathbb E|V_N|^4
                 -\frac5{96}\mathbb E|V_N|^6\right)
 \le t\le
 R\left(1-\frac{e_V}2+\frac3{16}\mathbb E|V_N|^4\right).
\]
Using the uniform moment bounds at $N=32$, the rational replay proves
\[
 \frac{347372732156351}{409600000000000}
       <t<\frac{2178061487133}{2560000000000},
 \qquad \ell<t<0.851\qquad(N\ge32).
\]
Consequently every one of these actual responses satisfies
\[
 \mathcal N<0.713254\qquad(N\ge32).
\]
For $N\ge341$ they also satisfy the cumulant-only region
\eqref{eq:cr-region}. The influence form explains why $N\ge32$ already
suffices on the exact support face.

There is a neighborhood statement, not just a family of exact responses.
For every $N\ge60$, any original response on the central interval with
\[
 D_t(V_N;q)\le1/200000
\]
satisfies $\mathcal N<0.713250$. Here its own cubic $F$ can have arbitrary
additional remote and visible components. Only the fixed test $V_N$
is a Wick sum. The proof substitutes
$\chi(V_N)\le81/(250\cdot60)$ in \eqref{eq:cr-main}.

\subsubsection{What this localizes, and what remains [P,?]}

Let
\[
 \eta_{\rm central}=\frac{240297199}{2651182000000}
 =\frac{500}{701}-\frac{2697232101}{3782000000}.
\]
Every actual central-row response with $\mathcal N\ge500/701$ must obey,
for \emph{every} remote cubic test,
\[
 2\sqrt{D_t(v;q)}+\chi(v)
       >\sqrt{\eta_{\rm central}/\varepsilon}.
\]
In particular, when the right side exceeds $\chi(v)$, the response must
be a definite distance inside that support constraint:
\[
 D_t(v;q)>\frac14
     \left(\sqrt{\eta_{\rm central}/\varepsilon}-\chi(v)\right)_+^2.
\]
This converts the remaining search into an explicit alternative:
large test cumulant/influence, or appreciable support gap. It replaces
a prescribed tensor frame by a criterion valid for every cubic test.

It does \emph{not} prove that all potentially extremizing tensors admit
a sufficiently good test. For example, at $q=0$ every test has
$D_t(v;0)\ge R-t$, and $v=0$ minimizes the present charge. The criterion
alone is then too weak to close the central row. Existing complementary
certificates or another argument are still needed there and on the rest
of the uncovered region. A full same-$\varepsilon$ row cover and the
Gaussian-to-matrix transfer remain necessary before claiming a larger
Grothendieck lower bound. The global certified interval is unchanged.

\paragraph{Replay and dependencies.}
The standard-library script \texttt{verify\_all.py} checks every
numerical constant above by rational arithmetic, including all global
derivative bounds, root enclosures, the full central interval, and the
infinite family $N\ge32$. It performs no quadrature and has no
truncation or rounding tail. The Gaussian proofs are the analytic
arguments in this section; finite arithmetic checks are not a substitute
for them. The remote input is still $C'$ from the manuscript, with its
sixth-moment assumption $C_6=4045313/100000$. The existing H1 and
influence inputs are used with their established finite-dimensional
Gaussian scope. No additional conjecture is assumed.

\paragraph{Small-$t$ corollary (IDE note, 21 October).} With $\rho\le\frac{96951}{100000}$ (\autoref{ch:spherical-kernel-cut}) and $\varepsilon=1891/2500$:
$0\le\|Ph\|\le\frac{239}{5000}$ implies objective $\le\frac{17831568694291}{25000000000000}<\frac{500}{701}$. The proof uses the exact objective identity and the duality
$\|Qk\|\le\rho$ for $|k|\le1$; no stationarity and no shared-sector restriction.

